% 10-federation.tex: independent authorities.
\section{Federation: independent authorities that recognise one another}
\label{sec:federation}

\subsection{Why there is no centre}

Two shapes were possible. A central instance holds every person's identity and every agency's
tokens under one trust root, with agencies as tenants. A federated shape gives each issuing
authority its own instance and its own signing key, records cross-authority trust explicitly, and
lets a relying party verify against published keys with no central service in the path. The
decision record chooses the second and states that it is not a free choice: the constitution names
the federation trust graph as load-bearing, says no agency holds a monopoly, and refuses
population-scale aggregation structurally. A central instance is a monopoly by construction and a
single store to seize, subpoena or mine. The federated shape is deliberately the more expensive one;
its cost is the inter-authority protocol, and the cost is accepted. \sImpl

\subsection{Trust is a row, not a chain}

An attestation is one row of \entity{AgencyTrustAttestation}: \emph{verifying agency V accepts
issuing agency I in context C, until a date}. Verification across agencies succeeds if and only if
exactly one active row matches the verifier, the issuer and the context. The lookup is a single
non-recursive query. It never computes a closure, so V trusting I never implies V trusts what I
trusts; cycles cannot arise because nothing walks the graph; and an operator who wants multi-hop
trust declares every hop, so the audit trail holds every decision rather than the inputs to an
inference. Self-attestation is refused by a check constraint. Revocation sets a date and a reason
together, once, and looks forward: events recorded while the edge was active are not invalidated.
\sImpl

One limit is named rather than glossed. The row records which Polaris operator created the
attestation; it is not a cryptographic signature from the attesting agency. Agency-level signing of
attestations is a migration the schema is shaped for and a key-management problem the readiness
ledger assigns to a deploying organisation. \sExt

\figfile{federation}{Two authorities and one directional attestation. B attests A's key in one
context; a relying party that trusts B accepts A's credential in that context, offline, from B's
manifest and A's revocation feed. A also trusts C, and that gives B nothing: trust is never a
closure. In CI this runs as two independent instances over HTTP under a mixed pair of parameter
sets.}{fig:federation}

\subsection{The three signed objects an authority publishes}

A relying party needs three things from a foreign authority, and all three are signed views over
append-only tables, published at short-lived endpoints and consumed offline by the detached
verifier. The federation manifest carries the authority's own anchors, every key it has ever used
with its real status, and the attestations it has made, each with the attested key. The epoch
checkpoint is its commitment to the latest point on its epoch chain. The revocation feed is the set
of revoked leaves. A manifest must be self-consistent, signed by one of the active anchors it
declares, so a stranger's key cannot sign one; authentic, under two witnesses; fresh, within a window
the consumer caps; and trusted, from an authority whose anchor the relying party chose. The
cross-authority decision then reads: accept if and only if the credential's signature is genuine,
the trusted manifest attests to that signing key in the presented context, and the issuer's feed,
fail-closed, does not list the credential. \sImpl

A holder can also prove membership in a foreign authority's epoch in zero knowledge. The proof
already carries the epoch root as a public input; the root already travels the federation signed in
the checkpoint. The decision composes the two: trust the foreign checkpoint through the graph to
obtain a trusted root, require the proof to bind to that root, epoch and context, then check the
proof by shelling to the prover binary as a local subprocess, with no network. If the binary is
absent the verifier abstains, which is never a false accept. The verdict names an authority and a
decision, never a token. \sImpl

\subsection{Two instances, and what they prove}

A dedicated CI job boots two independent instances as two authorities, each with its own database
and its own real ML-DSA root, talking only over HTTP, and drives the whole matrix: cross-verification
from a fetched manifest, attestation revocation flipping the decision, epoch checkpoints and
revocation feeds fetched and verified over the wire, a key compromise recorded on one side reversing
the other's acceptance, the exchange gateway, the timestamp service read out of a fetched registry,
the broker's login, document signing through the wallet, and a mixed pair of parameter sets
(ML-DSA-87 on one side, ML-DSA-65 on the other). It is red on any wrong decision. \sImpl

What the job proves is that the protocol works between two instances of this code base run by one
author on one machine. It does not prove institutional independence, and it does not prove that an
independent implementation, written by someone else from the specification, interoperates. The
first requires deployment by institutions that are not the author. The second requires an external
team, and has not happened. The repository draws this line itself: protocol independence,
independently operated institutions and independent third-party implementations are three different
claims, and only the first is demonstrated. \sExt

\begin{layercard}{Layer card: federation}
\qa{What it is}{Per-authority roots, directional per-context attestations, the signed manifest, epoch checkpoint and revocation feed, the aggregate status bundle, cross-authority zero-knowledge membership, and the two-instance drill.}
\qa{Why it exists}{No agency may hold a monopoly, and no single store of every person may exist; yet a relying party must be able to accept a credential issued elsewhere.}
\qa{What it trusts}{The anchors the relying party itself chose; the append-only tables the signed objects are views over.}
\qa{What it refuses}{Transitive trust; a central service in the verification path; a manifest not signed by its own declared anchor; an aggregator's word for a member's status; a callback the issuer could log.}
\qa{What it can see}{Published trust data: keys, attestations, epoch roots, revoked leaves. No person.}
\qa{What it cannot see}{The active population of any authority; who verified whom.}
\qa{If it fails}{A forged manifest fails self-consistency; a forked epoch chain or a rolled-back feed is non-repudiable evidence; a compromised issuer key is retired from an instant by a trust list the issuer does not control.}
\qa{Checked by}{\code{check\_federation\_topology}, \code{check\_inter\_authority\_protocol}, \code{check\_epoch\_revocation\_propagation}, \code{check\_cross\_authority\_zk}, \code{check\_federation\_two\_instances}, \code{check\_exchange\_trust\_directional}.}
\qa{Now or later}{Now, as a protocol between instances. Independent institutions and implementations are later, and are not ours to supply.}
\qa{Inside and outside}{Inside: use. Outside: transparency, because a signed history can still be told differently to different people.}
\end{layercard}

\nextlayer{Signed federation objects prove what an authority said. A dishonest authority can still
say different things to different observers, publishing one epoch chain to an auditor and another to
a relying party. Polaris therefore needs transparency: a way for observers to prove that history only
ever grew, and to compare notes.}
